The 1-Wire master (\peripheral{OWx}) is a Dallas/Maxim 1-Wire link-layer controller for the cheap-sensor ecosystem (DS18B20-class thermometers, iButton / DS24xx ROM-ID chips). It runs the five microsecond-scale bus primitives in hardware --- reset+presence, write-bit, read-bit, write-byte, read-byte --- off a programmable time base, and leaves the higher-level ROM search and CRC-8 to firmware over those primitives. It is master-only and supports both standard and overdrive speeds. The whole engine runs on the free-running fabric clock (MCLK); the single open-drain DQ pin is sampled through a two-flop synchronizer, so no flip-flop is ever clocked by the bus pin. The main features are listed below:

\begin{itemize}
	\item Five hardware link-layer primitives selected by \register{OWxCMD}: reset+presence detect, write-bit, read-bit, write-byte and read-byte, each launched by a single command write
	\item Byte transfers are LSB-first (the 1-Wire convention): \register{OWxTX} is shifted out least-significant-bit first, and read-byte assembles into \register{OWxRX} the same way
	\item Standard and overdrive timing sets, selected by \register{OWODS} and latched at each transaction launch so a mid-flight speed change never glitches a running slot
	\item A programmable time base (\register{OWxDIV}): one tick is $(\register{OWDIV}+1)$ MCLK cycles, giving a 0.5\,\textmu s tick at $\register{OWDIV}=11$ with the nominal 24\,MHz MCLK (all slot counts are integer half-microsecond ticks)
	\item Side-effect-free reads: the received byte/bit lands in \register{OWxRX}, never in the command register; a read never launches or clears anything
	\item Sticky write-1-to-clear status (\register{OWTCIF}, \register{OWNOPRES}, \register{OWSHORT}) and a single combined interrupt (vector 117)
	\item ROM search, ROM match and CRC-8 are done in firmware over the hardware bit primitives; the strong-pullup enable is a reserved stub (no driven-high phase this stage)
\end{itemize}

\subsection{The Time Base and the Two Speeds}

The slot state machine counts \emph{ticks}, never raw clock cycles, so all slot geometry is expressed in half-microsecond ticks independent of the actual MCLK rate. \register{OWxDIV} sets the tick period to $(\register{OWDIV}+1)$ MCLK cycles; program $\register{OWDIV}=11$ for the nominal 0.5\,\textmu s tick at 24\,MHz. \register{OWODS} selects between the standard slot-timing set and the overdrive (roughly ten-times-faster) set; it is latched into the transaction descriptor at launch, so writing a new \register{OWxCR} mid-transaction cannot disturb a slot already in flight. Because the engine clocks off the free-running fabric clock, its tick base is immune to clock reconfiguration, and software must \emph{not} apply the ``write \register{\register{SYSCLKCR}} to 0 first'' rule that the SMCLK peripherals need. A simulation bench uses a small \register{OWxDIV} to compress the microsecond scale into a few clock cycles per tick.

\subsection{Launching a Transaction}

A byte-lane-0 write to \register{OWxCMD} launches a transaction: it captures \{operation, bit-value, current speed, current \register{OWxTX} byte\} into a launch descriptor and starts the slot machine. The launch is suppressed when the master is disabled (\register{OWEN}$=0$) or already busy --- in that case the command content is still captured but no bus activity or completion occurs. Writes to \register{OWxTX}, \register{OWxCR} and \register{OWxDIV} never launch. \register{OWBUSY} covers the launch instant: it reads 1 in the same cycle as the launching command write, so firmware may poll \register{OWBUSY} to 0 immediately after writing \register{OWxCMD} with no blind window. Only one transaction may be outstanding --- poll \register{OWBUSY} clear before issuing the next command.

\subsection{Presence, Errors, and the Interrupt}

A reset transaction drives DQ low for the reset-low window, releases it, and samples the presence pulse: \register{OWPRES} sets if a device answered. \register{OWNOPRES} sets when a clean reset saw no presence pulse. \register{OWSHORT} sets when DQ is still low at the end of a recovery window --- a stuck or shorted bus --- after the master has released it. On a permanently-low bus during a reset, \emph{SHORT wins}: \register{OWSHORT} sets and \register{OWNOPRES} is suppressed, because presence is unevaluable on a bus that never returns high. The three event flags are sticky and write-1-to-clear (write a 1 to a bit to clear it; a read never clears them). The combined interrupt is $(\register{OWTCIF} \wedge \register{OWTCIE}) \vee ((\register{OWNOPRES} \vee \register{OWSHORT}) \wedge \register{OWERRIE})$, combinational and never latched; the interrupt service routine reads \register{OWxSR}, demultiplexes transaction-complete versus the error flags, and clears whichever fired before returning.

\subsection{Shared Access on \AsicNameForUserGuide}

In configurations that include it, \peripheral{OW0} lives in the shared peripheral window behind the multi-core arbiter (see Section \ref{s:multicore}) in the mutex-bank page, sub-slot 7 at \texttt{0x6700}, so any hart can use it. Its single combined interrupt source occupies vector 117 (transaction complete or error), delivered through the interrupt router like every other shared-peripheral source. This source sits above the external-interrupt (meip) delivery vector, which stays fixed at vector 85, above the I3C, NFC, \peripheral{GPIOx} port-4/port-5, RTC and PWM sources at vectors 86 through 116.

The block has one pin, the open-drain DQ line, on the alternate-function-2 plane of \texttt{P4.7} (\texttt{GPIO31}, the digital test port 3 pad) with its pull enabled at reset. That plane slot is one of the alternate-function output-spread copies of the TIMER0 compare-1 output, and DQ simply takes it over: compare 1 keeps its own primary pin \texttt{P3.1}, its two alternate-function-1 relocations, and more than twenty further spread copies, so nothing is lost. When \peripheral{OW0} is not instantiated the slot keeps its compare-1 copy and \texttt{P4.7} behaves exactly as before. Software routes DQ by selecting alternate function 2 for that pin's field in the port's \register{PxAFS} register; until it does, the receiver reads an idle-high line and the pad remains a plain digital test port.
